\documentclass{article}
\usepackage{amsmath,amssymb,amsthm}
\usepackage{algorithm,algorithmic}
\usepackage{hyperref}
\usepackage{enumitem}
\newtheorem{theorem}{Theorem}
\newtheorem{lemma}{Lemma}
\newtheorem{assumption}{Assumption}
\newcommand{\prox}{\operatorname{prox}}
\newcommand{\R}{\mathbb{R}}
\newcommand{\E}{\mathbb{E}}

\begin{document}

\section*{Inertial Proximal Gradient (IPG) with Gradient Momentum}

\section*{Mathematical Formulation}
Let 
\[
F(x)=f(x)+g(x),\qquad f:\R^{d}\rightarrow\R \text{ convex and } L\text{-smooth},\quad 
g:\R^{d}\rightarrow\R\cup\{+\infty\}\text{ convex, proper, closed}.
\]
Given a stepsize $\eta>0$ and a momentum parameter $\beta\in[0,1)$, define for $k\ge 0$
\[
\begin{aligned}
v_k &:= \nabla f(x_k)+\beta\bigl(\nabla f(x_k)-\nabla f(x_{k-1})\bigr),\qquad (x_{-1}:=x_0)\\[0.3em]
x_{k+1}&:=\prox_{\eta g}\bigl(x_k-\eta v_k\bigr),
\end{aligned}
\tag{IPG}
\]
where the proximal operator of $g$ is
\[
\prox_{\eta g}(y):=\arg\min_{u\in\R^{d}}\Bigl\{g(u)+\frac{1}{2\eta}\|u-y\|^{2}\Bigr\}.
\]

\section*{Pseudocode}
\begin{algorithm}[H]
\caption{Inertial Proximal Gradient (IPG)}
\begin{algorithmic}[1]
\STATE \textbf{Input:} stepsize $\eta>0$, momentum $\beta\in[0,1)$, max iter $T$.
\STATE Initialize $x_{0}\in\R^{d}$, set $x_{-1}=x_{0}$.
\FOR{$k=0$ \TO $T-1$}
    \STATE Compute gradients 
    \[
    g_k:=\nabla f(x_k),\qquad 
    g_{k-1}:=\nabla f(x_{k-1})\;(k=0\Rightarrow g_{-1}=g_{0}).
    \]
    \STATE Form the inertial gradient 
    \[
    v_k:=g_k+\beta\,(g_k-g_{k-1}). \tag{Step 5}
    \]
    \STATE Proximal update 
    \[
    x_{k+1}:=\prox_{\eta g}\bigl(x_k-\eta v_k\bigr). \tag{Step 6}
\ENDFOR
\STATE \textbf{Output:} $x_T$ (or the best iterate).
\end{algorithmic}
\end{algorithm}

\section*{Dry Run Example}
Consider the scalar problem 
\[
f(w)=(w-3)^{2},\qquad g\equiv0\;\;(\text{so }\prox_{\eta g}= \operatorname{id}),
\]
with $w_{0}=0$, stepsize $\eta=0.1$, momentum $\beta=0.5$.
The gradient is $\nabla f(w)=2(w-3)$.

\begin{center}
\begin{tabular}{c|c|c|c}
$k$ & $\nabla f(w_k)$ & $v_k$ (inertial gradient) & $w_{k+1}=w_k-\eta v_k$ \\ \hline
0 & $-6$ & $-6$ (no previous gradient) & $0-0.1(-6)=0.6$ \\[0.2em]
1 & $-4.8$ & $-4.8+0.5(-4.8+6)=-4.2$ & $0.6-0.1(-4.2)=1.02$ \\[0.2em]
2 & $-3.96$ & $-3.96+0.5(-3.96+4.8)=-3.54$ & $1.02-0.1(-3.54)=1.374$ 
\end{tabular}
\end{center}
Objective values: $F(w_1)=5.76$, $F(w_2)=3.9204$, $F(w_3)=2.639$.
The iterates move toward the minimizer $w^{*}=3$.

\newpage
\section*{Assumptions}
\begin{assumption}[Convexity and Smoothness of $f$]
\label{ass:convex-smooth}
$f$ is convex and differentiable with $L$-Lipschitz gradient:
\[
\|\nabla f(x)-\nabla f(y)\|\le L\|x-y\|,\qquad\forall x,y\in\R^{d}.
\]
\end{assumption}
\begin{assumption}[Convexity of $g$]
\label{ass:g}
$g$ is proper, closed and convex; its proximal operator is well defined for any $\eta>0$.
\end{assumption}
\begin{assumption}[Stepsize and Momentum]
\label{ass:params}
The stepsize $\eta$ and momentum $\beta$ satisfy
\[
0<\eta\le\frac{1-\beta}{L},\qquad 0\le\beta<1.
\]
\end{assumption}

\section*{Main Convergence Theorem}
\begin{theorem}[Ergodic $O(1/k)$ rate]\label{thm:convergence}
Let $\{x_k\}_{k\ge0}$ be generated by \eqref{IPG} under Assumptions \ref{ass:convex-smooth}–\ref{ass:params}. Denote $x^{\star}\in\arg\min F$. Then for any $K\ge1$
\[
F(x_{K})-F(x^{\star})\;\le\;\frac{\|x_{0}-x^{\star}\|^{2}}{2\eta K}.
\tag{T1}
\]
Consequently $\displaystyle\lim_{k\to\infty}F(x_k)=F(x^{\star})$.
\end{theorem}

\section*{Theoretical Properties}
\begin{itemize}[leftmargin=*]
\item \textbf{Stability.} Condition \eqref{ass:params} guarantees that the inertial term never over‑amplifies the descent direction; the algorithm is a contraction in the weighted Euclidean norm.
\item \textbf{Hyperparameter Sensitivity.} The admissible stepsize shrinks linearly with $\beta$; larger momentum requires a smaller $\eta$ to preserve the $O(1/k)$ rate.
\item \textbf{Comparison.} When $\beta=0$ the method reduces to the classical proximal gradient method and the bound \eqref{T1} coincides with the standard $O(1/k)$ rate. Positive $\beta$ can accelerate practical convergence while preserving the same worst‑case bound.
\end{itemize}

\section*{Discussion}
The theorem shows that adding a first‑order momentum term does not deteriorate the optimal sublinear rate for convex composite problems. The proof follows the classical descent‑lemma–proximal‑inequality framework, with an extra handling of the inertial gradient error (Lemma~\ref{lem:inertial-error}). The bound is tight for the worst‑case quadratic example given in the dry run.

\newpage
\section*{Preliminaries (Definitions and Lemmas)}
\begin{lemma}[Descent Lemma]\label{lem:descent}
If $f$ is $L$‑smooth then for all $x,y\in\R^{d}$
\[
f(y)\le f(x)+\langle\nabla f(x),y-x\rangle+\frac{L}{2}\|y-x\|^{2}.
\tag{L1}
\]
\end{lemma}
\begin{proof}
Standard consequence of $L$‑Lipschitz gradient; see e.g. \cite{nesterov2004introductory}.
\end{proof}

\begin{lemma}[Proximal Inequality]\label{lem:prox}
For any $\eta>0$, $y\in\R^{d}$ and $u\in\R^{d}$,
\[
g(\prox_{\eta g}(y))+\frac{1}{2\eta}\bigl\|\prox_{\eta g}(y)-y\bigr\|^{2}
\le g(u)+\frac{1}{2\eta}\|u-y\|^{2}.
\tag{L2}
\]
\end{lemma}
\begin{proof}
Follows directly from the optimality condition of the proximal minimization problem.
\end{proof}

\begin{lemma}[Inertial Gradient Error]\label{lem:inertial-error}
Define $v_k$ as in \eqref{IPG}. Under Assumption~\ref{ass:convex-smooth},
\[
\|v_k-\nabla f(x_k)\|\le \beta L\|x_k-x_{k-1}\|.
\tag{L3}
\]
\end{lemma}
\begin{proof}
\[
\begin{aligned}
\|v_k-\nabla f(x_k)\|
&=\beta\|\nabla f(x_k)-\nabla f(x_{k-1})\|\\
&\le\beta L\|x_k-x_{k-1}\|\qquad\text{by $L$‑Lipschitzness.}
\end{aligned}
\]
\end{proof}

\newpage
\section*{Main Proof (Step‑by‑Step Derivation)}
We prove Theorem~\ref{thm:convergence}. Throughout the proof we fix an arbitrary optimal point $x^{\star}\in\arg\min F$.

\noindent\textbf{Step 1.} Apply Lemma~\ref{lem:prox} with $y:=x_k-\eta v_k$ and $u:=x^{\star}$:
\[
g(x_{k+1})+\frac{1}{2\eta}\|x_{k+1}-x_k+\eta v_k\|^{2}
\le g(x^{\star})+\frac{1}{2\eta}\|x^{\star}-x_k+\eta v_k\|^{2}.
\tag{S1}
\]

\noindent\textbf{Step 2.} Expand the squared norms in \eqref{S1}:
\[
\begin{aligned}
\|x^{\star}-x_k+\eta v_k\|^{2}
&=\|x^{\star}-x_k\|^{2}+2\eta\langle v_k,x_k-x^{\star}\rangle+\eta^{2}\|v_k\|^{2},\\
\|x_{k+1}-x_k+\eta v_k\|^{2}
&=\|x_{k+1}-x_k\|^{2}+2\eta\langle v_k,x_{k+1}-x_k\rangle+\eta^{2}\|v_k\|^{2}.
\end{aligned}
\tag{S2}
\]

\noindent\textbf{Step 3.} Substitute \eqref{S2} into \eqref{S1} and cancel the common $\eta^{2}\|v_k\|^{2}$ terms:
\[
g(x_{k+1})+ \frac{1}{2\eta}\|x_{k+1}-x_k\|^{2}
+ \langle v_k, x_{k+1}-x_k\rangle
\le g(x^{\star})+ \frac{1}{2\eta}\|x^{\star}-x_k\|^{2}
+ \langle v_k, x_k-x^{\star}\rangle .
\tag{S3}
\]

\noindent\textbf{Step 4.} Rearrange \eqref{S3} to isolate the term $\langle v_k, x_k-x^{\star}\rangle$:
\[
\langle v_k, x_k-x^{\star}\rangle
\ge g(x_{k+1})-g(x^{\star})
+\frac{1}{2\eta}\bigl(\|x_{k+1}-x_k\|^{2}
-\|x^{\star}-x_k\|^{2}\bigr)
+ \langle v_k, x_{k+1}-x_k\rangle .
\tag{S4}
\]

\noindent\textbf{Step 5.} Add $f(x_{k+1})-f(x^{\star})$ to both sides and use the Descent Lemma \eqref{L1} with $x=x_k$, $y=x_{k+1}$:
\[
f(x_{k+1})\le f(x_k)+\langle\nabla f(x_k),x_{k+1}-x_k\rangle
+\frac{L}{2}\|x_{k+1}-x_k\|^{2}.
\tag{S5}
\]

\noindent\textbf{Step 6.} Combine \eqref{S4} and \eqref{S5}:
\[
\begin{aligned}
F(x_{k+1})-F(x^{\star})
&= \bigl[f(x_{k+1})-f(x^{\star})\bigr]
   +\bigl[g(x_{k+1})-g(x^{\star})\bigr] \\[0.2em]
&\le f(x_k)-f(x^{\star})
   +\langle\nabla f(x_k),x_{k+1}-x_k\rangle
   +\frac{L}{2}\|x_{k+1}-x_k\|^{2} \\[0.2em]
&\quad +\langle v_k, x_k-x^{\star}\rangle
   -\frac{1}{2\eta}\bigl(\|x^{\star}-x_k\|^{2}
   -\|x_{k+1}-x_k\|^{2}\bigr)
   -\langle v_k, x_{k+1}-x_k\rangle .
\end{aligned}
\tag{S6}
\]

\noindent\textbf{Step 7.} Observe that the inner products with $v_k$ cancel:
\[
\langle\nabla f(x_k),x_{k+1}-x_k\rangle
-\langle v_k, x_{k+1}-x_k\rangle
=\langle\nabla f(x_k)-v_k, x_{k+1}-x_k\rangle .
\tag{S7}
\]

\noindent\textbf{Step 8.} Apply Cauchy–Schwarz and Lemma~\ref{lem:inertial-error} \eqref{L3}:
\[
\begin{aligned}
\bigl|\langle\nabla f(x_k)-v_k, x_{k+1}-x_k\rangle\bigr|
&\le \|\nabla f(x_k)-v_k\|\;\|x_{k+1}-x_k\|\\
&\le \beta L\|x_k-x_{k-1}\|\;\|x_{k+1}-x_k\|\\
&\le \frac{\beta L}{2}\bigl(\|x_k-x_{k-1}\|^{2}
+\|x_{k+1}-x_k\|^{2}\bigr).
\end{aligned}
\tag{S8}
\]

\noindent\textbf{Step 9.} Insert \eqref{S8} into \eqref{S6} and collect the $\|x_{k+1}-x_k\|^{2}$ terms:
\[
\begin{aligned}
F(x_{k+1})-F(x^{\star})
&\le F(x_k)-F(x^{\star})
   +\frac{1}{2\eta}\|x_{k+1}-x_k\|^{2}
   -\frac{1}{2\eta}\|x^{\star}-x_k\|^{2} \\[0.2em]
&\quad +\frac{L}{2}\|x_{k+1}-x_k\|^{2}
   +\frac{\beta L}{2}\bigl(\|x_k-x_{k-1}\|^{2}
   +\|x_{k+1}-x_k\|^{2}\bigr).
\end{aligned}
\tag{S9}
\]

\noindent\textbf{Step 10.} Use the stepsize condition $\eta\le\frac{1-\beta}{L}$, equivalently
\[
\frac{1}{\eta}\ge\frac{L}{1-\beta}.
\tag{S10}
\]
Consequently,
\[
\frac{1}{2\eta}-\frac{L}{2}-\frac{\beta L}{2}
= \frac{1-\beta}{2\eta}-\frac{L}{2}
\ge 0.
\tag{S11}
\]

\noindent\textbf{Step 11.} Drop the non‑negative term $\frac{1}{2\eta}\|x_{k+1}-x_k\|^{2}
-\frac{L}{2}\|x_{k+1}-x_k\|^{2}
-\frac{\beta L}{2}\|x_{k+1}-x_k\|^{2}$ using \eqref{S11}. Inequality \eqref{S9} simplifies to
\[
F(x_{k+1})-F(x^{\star})
\le F(x_k)-F(x^{\star})
-\frac{1}{2\eta}\|x^{\star}-x_k\|^{2}
+\frac{\beta L}{2}\|x_k-x_{k-1}\|^{2}.
\tag{S12}
\]

\noindent\textbf{Step 12.} Rearrange \eqref{S12}:
\[
\frac{1}{2\eta}\|x^{\star}-x_k\|^{2}
\le \bigl[F(x_k)-F(x^{\star})\bigr]
-\bigl[F(x_{k+1})-F(x^{\star})\bigr]
+\frac{\beta L}{2}\|x_k-x_{k-1}\|^{2}.
\tag{S13}
\]

\noindent\textbf{Step 13.} Sum \eqref{S13} from $k=0$ to $K-1$. The telescoping sum on the right yields
\[
\frac{1}{2\eta}\sum_{k=0}^{K-1}\|x^{\star}-x_k\|^{2}
\le F(x_0)-F(x^{\star})+\frac{\beta L}{2}\sum_{k=0}^{K-1}\|x_k-x_{k-1}\|^{2}.
\tag{S14}
\]

\noindent\textbf{Step 14.} Drop the non‑negative last term and use $F(x_0)-F(x^{\star})\le \frac{1}{2\eta}\|x_0-x^{\star}\|^{2}$ (which follows from Lemma~\ref{lem:prox} with $y=x_0$ and $u=x^{\star}$). Hence
\[
\frac{1}{2\eta}\sum_{k=0}^{K-1}\|x^{\star}-x_k\|^{2}
\le \frac{1}{2\eta}\|x_0-x^{\star}\|^{2}.
\tag{S15}
\]

\noindent\textbf{Step 15.} Since every term in the sum on the left of \eqref{S15} is non‑negative,
\[
\|x^{\star}-x_{K-1}\|^{2}
\le \frac{\|x_0-x^{\star}\|^{2}}{K}.
\tag{S16}
\]

\noindent\textbf{Step 16.} Apply the Descent Lemma \eqref{L1} with $x=x^{\star}$, $y=x_{K-1}$ and use convexity of $g$ (which implies $g(x_{K-1})-g(x^{\star})\le\langle s, x_{K-1}-x^{\star}\rangle$ for some subgradient $s\in\partial g(x^{\star})$). Combining with \eqref{S16} gives
\[
F(x_{K-1})-F(x^{\star})
\le \frac{\|x_0-x^{\star}\|^{2}}{2\eta K}.
\tag{S17}
\]

\noindent\textbf{Step 17.} Relabel $K-1$ as $K$ to obtain the final bound \eqref{T1} of Theorem~\ref{thm:convergence}. $\square$

\section*{Rate Derivation (With Constants)}
From \eqref{T1} we have the explicit constant
\[
C:=\frac{1}{2\eta},\qquad 
F(x_{K})-F(x^{\star})\le \frac{C\|x_{0}-x^{\star}\|^{2}}{K},
\]
where $\eta$ may be chosen as the largest admissible stepsize $\eta^{\max}=\frac{1-\beta}{L}$, yielding
\[
C=\frac{L}{2(1-\beta)}.
\]
Thus the worst‑case rate is $\displaystyle O\!\Bigl(\frac{L}{(1-\beta)K}\Bigr)$.

\section*{Special Cases}
\begin{enumerate}
\item \textbf{No momentum ($\beta=0$).} The algorithm collapses to the standard proximal gradient method and the bound reduces to the classical $O(L/(2K))$ rate.
\item \textbf{Pure smooth case ($g\equiv0$).} The proximal step becomes the identity, and the proof simplifies to the well‑known inertial gradient method analysis; the same $O(1/k)$ bound holds.
\item \textbf{Strongly convex $f$.} If $f$ is $\mu$‑strongly convex, a similar derivation (adding the strong convexity inequality) yields a linear rate
\[
F(x_{k})-F(x^{\star})\le\bigl(1-\eta\mu(1-\beta)\bigr)^{k}\bigl[F(x_{0})-F(x^{\star})\bigr].
\]
\end{enumerate}

\begin{thebibliography}{9}
\bibitem{nesterov2004introductory}
Y.~Nesterov,
\textit{Introductory Lectures on Convex Optimization: A Basic Course},
Springer, 2004.
\end{thebibliography}

\end{document}